\section*{Bab \ref{ch:g12:matrix} --- Matriks dan Graf}

\begin{solution}{exo:g12:matrix:1}
\[
A + B = \begin{pmatrix} 3 & 2\\ 1 & 2\end{pmatrix}, \quad
AB = \begin{pmatrix} 4 & 2\\ 1 & 1\end{pmatrix}, \quad
BA = \begin{pmatrix} 2 & 4\\ 1 & 3\end{pmatrix}, \quad
A^2 = \begin{pmatrix} 1 & 4\\ 0 & 1\end{pmatrix}.
\]
Perhatikan bahwa $AB \neq BA$.
\end{solution}

\begin{solution}{exo:g12:matrix:2}
$\det A = 6 - 5 = 1 \neq 0$:
$A^{-1} = \begin{pmatrix} 3 & -5\\ -1 & 2 \end{pmatrix}$.
$\det B = 12 - 12 = 0$: jadi $B$ tak \omterm{def:g12:matrix:inverse}{dapat dibalik}.
\end{solution}

\begin{solution}{exo:g12:matrix:3}
Sistemnya adalah $AX = Y$ dengan $A$ seperti pada \cref{exo:g12:matrix:2} dan
$Y = \begin{pmatrix} 1\\ 2\end{pmatrix}$:
\[
X = A^{-1}Y = \begin{pmatrix} 3 & -5\\ -1 & 2\end{pmatrix}
\begin{pmatrix} 1\\ 2\end{pmatrix}
= \begin{pmatrix} -7\\ 3\end{pmatrix}:
\qquad x = -7,\ y = 3 .
\]
\end{solution}

\begin{solution}{exo:g12:matrix:4}
\emph{1.} $N^2 = \begin{pmatrix} 0&1\\0&0\end{pmatrix}
\begin{pmatrix} 0&1\\0&0\end{pmatrix} = 0$, dan $I$ komutatif dengan setiap
\omterm{def:g12:matrix:matrix}{matriks}.

\emph{2.} Karena kedua sukunya komutatif, teorema binomial berlaku dan semua
suku yang memuat $N^2$ lenyap:
\[
A^k = (2I + N)^k = 2^k I + k\,2^{k-1} N
= \begin{pmatrix} 2^k & k\,2^{k-1}\\ 0 & 2^k \end{pmatrix}.
\]
Periksalah untuk $k = 2$:
$A^2 = \begin{pmatrix} 2&1\\0&2\end{pmatrix}^2
= \begin{pmatrix} 4&4\\0&4\end{pmatrix}$, dan rumusnya memberi
$2^2 = 4$, $2 \times 2 = 4$. \checkmark
\end{solution}

\begin{solution}{exo:g12:matrix:5}
\emph{Induksinya.} Untuk $n = 1$:
$A^1 = \begin{pmatrix} 0&1\\1&1\end{pmatrix}
= \begin{pmatrix} F_0 & F_1\\ F_1 & F_2\end{pmatrix}$. Andaikan rumusnya benar
untuk $n$; maka
\[
A^{n+1} = A^n A
= \begin{pmatrix} F_{n-1} & F_n\\ F_n & F_{n+1}\end{pmatrix}
\begin{pmatrix} 0 & 1\\ 1 & 1\end{pmatrix}
= \begin{pmatrix} F_n & F_{n-1} + F_n\\ F_{n+1} & F_n + F_{n+1}\end{pmatrix}
= \begin{pmatrix} F_n & F_{n+1}\\ F_{n+1} & F_{n+2}\end{pmatrix}.
\]
\emph{Identitasnya.} Untuk \omterm{def:g12:matrix:matrix}{matriks} $2\times2$, penjabarannya menunjukkan
$\det(MN) = \det M \det N$; karena itu
$\det(A^n) = (\det A)^n = (-1)^n$, sedangkan
$\det A^n = F_{n-1}F_{n+1} - F_n^2$. (Inilah \emph{identitas Cassini}.)
\end{solution}

\begin{solution}{exo:g12:matrix:6}
\emph{1.} Dengan mengurutkan simpul $1, 2, 3$:
\[
M = \begin{pmatrix} 0&1&1\\ 0&0&1\\ 1&0&0\end{pmatrix}, \quad
M^2 = \begin{pmatrix} 1&0&1\\ 1&0&0\\ 0&1&1\end{pmatrix}, \quad
M^3 = \begin{pmatrix} 1&1&1\\ 1&0&1\\ 1&0&1 \end{pmatrix}.
\]

\emph{2.} $\bigl(M^3\bigr)_{11} = 1$: jadi tepat ada satu jalan tertutup sepanjang $3$
di simpul $1$, yaitu $1 \to 2 \to 3 \to 1$. (Jalan
$1 \to 3 \to 1$ hanya sepanjang $2$, sedangkan $1 \to 3$ lalu $3\to1$ lalu
$1\to3$ berakhir di $3$.)
\end{solution}

\begin{solution}{exo:g12:matrix:7}
\emph{1.} $a_{n+1} = 0.8a_n + 0.3b_n$, $b_{n+1} = 0.2a_n + 0.7b_n$, jadi
$M = \begin{pmatrix} 0.8 & 0.3\\ 0.2 & 0.7\end{pmatrix}$.

\emph{2.} $MX = X$ memberi $0.8a + 0.3b = a$, \emph{yaitu} $0.3b = 0.2a$, sehingga
$b = \frac23 a$; lalu dengan $a + b = 1$: $a = 0.6$, $b = 0.4$.

\emph{3.} Dengan memakai $b_n = 1 - a_n$:
$a_{n+1} = 0.8a_n + 0.3(1 - a_n) = 0.5a_n + 0.3$, sehingga
\[
c_{n+1} = a_{n+1} - 0.6 = 0.5a_n + 0.3 - 0.6 = 0.5(a_n - 0.6) = 0.5\,c_n .
\]
Jadi $c_n = 0.5^n c_0 \to 0$: maka $a_n \to 0.6$ dan $b_n \to 0.4$, berapa pun
\omterm{def:g11:prob:rv}{distribusi} awalnya.
\end{solution}

\begin{solution}{exo:g12:matrix:8}
\emph{1.} $\det P = -2$, jadi
$P^{-1} = -\frac12\begin{pmatrix} -1 & -1\\ -1 & 1\end{pmatrix}
= \frac12\begin{pmatrix} 1 & 1\\ 1 & -1\end{pmatrix}$. Maka
\[
AP = \begin{pmatrix} 4 & 2\\ 4 & -2 \end{pmatrix}, \qquad
D = P^{-1}AP = \frac12\begin{pmatrix} 1&1\\1&-1\end{pmatrix}
\begin{pmatrix} 4&2\\4&-2\end{pmatrix}
= \begin{pmatrix} 4 & 0\\ 0 & 2\end{pmatrix}.
\]

\emph{2.} Dari $A = PDP^{-1}$, induksi yang langsung memberi
$A^n = PD^nP^{-1}$ dengan
$D^n = \begin{pmatrix} 4^n & 0\\ 0 & 2^n\end{pmatrix}$, sehingga
\[
A^n = P D^n P^{-1}
= \begin{pmatrix} 4^n & 2^n\\ 4^n & -2^n\end{pmatrix}\cdot
\frac12\begin{pmatrix} 1&1\\1&-1\end{pmatrix}
= \frac12\begin{pmatrix} 4^n + 2^n & 4^n - 2^n\\
4^n - 2^n & 4^n + 2^n\end{pmatrix}.
\]
Menerapkan $A^n$ pada $X_0 = \begin{pmatrix} u_0\\v_0\end{pmatrix}$ menghasilkan
kembali persis rumus pada \cref{ex:g12:matrix:coupled}.
\end{solution}

\begin{solution}{pb:g12:matrix:1}
\textbf{1.} $AB = \begin{pmatrix} 2 & 1\\ 4 & 3\end{pmatrix}$
dan $BA = \begin{pmatrix} 3 & 4\\ 1 & 2\end{pmatrix}$: jadi perkalian
\omterm{def:g12:matrix:matrix}{matriks} tidak komutatif --- $B$ menukar kolom bila di
kanan, dan menukar baris bila di kiri.

\textbf{2.} \omterm{def:g11:vect:det}{Determinannya} $6 - 5 = 1$, jadi balikannya
$\begin{pmatrix} 3 & -1\\ -5 & 2\end{pmatrix}$. Dengan menerapkannya pada
$\binom{4}{7}$: $x = 12 - 7 = 5$, $y = -20 + 14 = -6$.

\textbf{3.} $N^2 = 0$. Maka $(I + N)^n = I + nN$ lewat induksi:
$(I + nN)(I + N) = I + (n+1)N + nN^2 = I + (n+1)N$.

\textbf{4.} $A = \begin{pmatrix} 0&1&1\\ 1&0&1\\ 1&1&0
\end{pmatrix}$, dan $A^3$ berisian diagonal $2$: jadi dari tiap
simpulnya ada tepat dua jalan tertutup sepanjang $3$ (segitiganya
ditempuh searah atau berlawanan arah jarum jam) --- teorema pencacahan
itu sedang beraksi.

\textbf{5.} $D^n = \begin{pmatrix} 2^n & 0\\ 0 & 2^{-n}
\end{pmatrix}$: satu arahnya meledak, arah lainnya mati ---
nasib diagonalnya adalah \omterm{def:g12:seq:sequence}{barisan} ukur yang saling bebas.

\textbf{6.} $F^2 = \begin{pmatrix} 2 & 1\\ 1 & 1\end{pmatrix}$,
$F^3 = \begin{pmatrix} 3 & 2\\ 2 & 1\end{pmatrix}$,
$F^4 = \begin{pmatrix} 5 & 3\\ 3 & 2\end{pmatrix}$: Fibonacci
di mana-mana; terkaannya seperti yang dinyatakan.

\textbf{7.} Jika $F^n = \begin{pmatrix} F_{n+1} & F_n\\ F_n &
F_{n-1}\end{pmatrix}$, maka
\[
F^{n+1} = F^n F =
\begin{pmatrix} F_{n+1} + F_n & F_{n+1}\\
F_n + F_{n-1} & F_n \end{pmatrix}
= \begin{pmatrix} F_{n+2} & F_{n+1}\\ F_{n+1} & F_n
\end{pmatrix} :
\]
itulah pewarisannya; sedangkan kasus dasarnya $n = 1$ adalah $F$ sendiri, dengan
memakai kesepakatan $F_0 = 0$ (yang memperluas rekurensinya ke belakang).

\textbf{8.} $\det F = -1$, jadi
$\det(F^n) = (\det F)^n = (-1)^n$; dan secara langsung
$\det(F^n) = F_{n+1}F_{n-1} - F_n^2$: itulah Cassini, dalam satu baris. (Kaidah
hasil kali bagi \omterm{def:g11:vect:det}{determinan} $2 \times 2$ adalah penjabaran lima menit
yang menyenangkan.)

\textbf{9.} Kanan atas $F^m F^n$:
$F_{m+1}F_n + F_m F_{n-1}$; kanan atas $F^{m+n}$:
$F_{m+n}$. Untuk $m = n = 3$:
$F_4 F_3 + F_3 F_2 = 3 \times 2 + 2 \times 1 = 8 = F_6$.

\textbf{10.} Untuk $k = 1$: sepele. Jika $F_n \mid F_{kn}$, maka
rumus penjumlahannya dengan $m = kn$ memberi
$F_{(k+1)n} = F_{kn+1}F_n + F_{kn}F_{n-1}$: kedua sukunya
kelipatan $F_n$. Jadi $F_n \mid F_{kn}$ untuk semua $k$: periksalah bahwa
$F_3 = 2$ \omterm{def:g12:arith:divides}{membagi} $F_6 = 8$ dan $F_9 = 34$.

\textbf{11.} $F^{100} = F^{64} F^{32} F^4$: yaitu tujuh kali pengkuadratan
($F^2, F^4, \dots, F^{64}$) ditambah dua penggabungan --- sembilan
kali perkalian, bukan sembilan puluh sembilan. Inilah
muslihat tabel pelipatduaan para juru tulis Mesir, yang diangkat dari
bilangan ke \omterm{def:g12:matrix:matrix}{matriks}: tulislah $100$ dalam biner, lalu kalikan
pelipatduaan yang kamu perlukan.

\textbf{12.} $0.8 + 0.2 = 1$ dan $0.4 + 0.6 = 1$: besok
haruslah \emph{suatu} cuaca --- tiap kolomnya adalah \omterm{def:g11:prob:rv}{distribusi}
peluang yang lengkap, sehingga peluangnya kekal.

\textbf{13.} Besok: $\binom{0.8}{0.2}$. Lusa:
$M\binom{0.8}{0.2} = \binom{0.72}{0.28}$.

\textbf{14.} $Mv = v$ dengan $v = \binom{s}{r}$, $s + r = 1$:
maka $0.8s + 0.4r = s$ memberi $0.4r = 0.2s$, jadi $s = 2r$ dan
$v = \binom{2/3}{1/3}$. Dalam jangka panjang, dua hari dari tiga hari
cerah --- bagaimanapun rupa hari ini.

\textbf{15.} Dari $\binom01$: komponen cerahnya $0.4$, $0.56$,
$0.624$, $0.6496$; jurangnya terhadap $\frac23$: $0.267$, $0.107$,
$0.043$, $0.017$ --- jadi tiap langkahnya mengalikan jurangnya tepat dengan
$0.4$ (yaitu nilai eigen kedua mesinnya): kekonvergenan
ukur menuju keadaan mapannya.

\textbf{16.} Kolomnya (dari $A$, $B$, $C$):
$P = \begin{pmatrix} 0 & 0 & 1\\ \frac12 & 0 & 0\\
\frac12 & 1 & 0\end{pmatrix}$. Keadaan mapannya: $v_A = v_C$,
$v_B = \frac{v_A}{2}$, $v_C = \frac{v_A}{2} + v_B$; lalu dengan berjumlah
$1$: $v = \left(\frac25, \frac15, \frac25\right)$. Peringkatnya:
$A$ dan $C$ seri di tempat pertama, $B$ terakhir.

\textbf{17.} Halaman $C$ menerima \emph{seluruh} lalu lintas $B$ dan separuh
milik $A$, lalu menyalurkan semuanya kembali ke $A$: keadaan
mapannya mengukur di mana sang peselancar \emph{menghabiskan waktunya}, bukan
berapa banyak tautan yang menunjuk masuk --- satu tautan dari halaman populer
mengalahkan beberapa tautan dari halaman sepi. Pembobotan rekursif itu
justru gagasan pendiri Google; peredamnya mengurus perangkap laba-laba
dan jalan buntu.

\textbf{18.} $A^n = P D^n P^{-1}$ memberi
$u_n = \frac{4^n + 2^n}{2}$ (dan
$v_n = \frac{4^n - 2^n}{2}$). Periksalah: $u_1 = 3$, $u_2 = 10$,
$u_3 = 36$; sedangkan secara langsung: $(1,0) \to (3,1) \to (10,6) \to
(36, 28)$, jadi cocok.

\textbf{19.} Pendiagonalan mengganti ke \omterm{def:g10:coordgeom:system}{koordinat} yang membuat
sistem berpasangan itu terurai menjadi \omterm{def:g12:seq:sequence}{barisan} ukur yang saling
bebas --- tiap nilai eigennya berlari dalam lombanya sendiri. Keadaan
mapan Markov adalah \omterm{def:g11:vect:vector}{vektor} eigen dengan nilai eigen $1$, dan
laju kekonvergenan pada pertanyaan 15 adalah nilai eigen berikutnya: jadi
mesin cuaca itu sejak awal sebuah kisah eigen.

\textbf{20.} Pembukuan: satu sistem adalah satu \omterm{def:g10:algebra:equation}{persamaan} \omterm{def:g12:matrix:matrix}{matriks},
yang diselesaikan oleh satu balikan. Kombinatorika: pangkat \omterm{def:g12:matrix:graph}{matriks
ketetanggaan} mencacah jalan, tautan, dan hubungan. Perubahan: pangkat
mesinnya membawa keadaan menuju nasibnya, dan
arah eigennya (yaitu arah emas milik Fibonacci, keadaan mapan cuaca,
\omterm{def:g11:vect:vector}{vektor} peringkat milik jejaring) adalah nasib itu sendiri.
Aljabar linear, di jilid universitas, adalah ilmu tentang
justru hal ini.

\end{solution}
